\SourcePage{723}{726}
\section{The scattering matrix in the presence of reactions}

The cross-section $\sigma_r$ considered in \S~142, 143 was the total
cross-section of all possible inelastic scattering channels. We shall now
show how to construct the general theory of inelastic collisions in which
each channel can be considered separately.

Suppose that a collision of two particles again produces two particles,
whether the same or different. Number all reaction channels possible at the
specified energy and attach the corresponding indices to the quantities
belonging to them.

Let channel $i$ be the entrance channel. In this channel the wave function of
the relative motion in the center-of-mass system is the familiar sum of an
incident plane wave and an elastically scattered outgoing wave:
\begin{equation}
 \psi_i=\exp(ik_i z)+f_{ii}(\theta)\frac{\exp(ik_i r)}r.
 \tag{144.1}
\end{equation}
The squared amplitude $f_{ii}$ gives the elastic scattering cross-section in
channel $i$:
\begin{equation}
 d\sigma_{ii}=|f_{ii}|^2d\Omega.
 \tag{144.2}
\end{equation}

In the other channels, indexed by $f$, the wave functions of relative motion
are outgoing waves. For a reason that will become clear below, it is
convenient to write them as\footnote{As before (compare the note on p.~187),
the initial state is denoted by $i$ and the final state by $f$. In the
scattering amplitude, the final-state index is placed to the left of the
initial-state index, in accordance with the order of matrix-element indices.
For uniformity, the indices in cross-section symbols will be ordered in the
same way.}
\begin{equation}
 \psi_f=f_{fi}(\theta)\sqrt{m_f/m_i}\,\frac{\exp(ik_f r)}r,
 \tag{144.3}
\end{equation}

\SourcePage{724}{727}
where $\mathbf{k}_f$ is the wave vector of relative motion of the reaction
products in channel $f$, $\theta$ is its angle with the $z$ axis, and $m_i$,
$m_f$ are the reduced masses of the two initial and two final particles. The
scattered current in the solid-angle element $d\Omega_f$ is obtained by
multiplying $|\psi_f|^2$ by $v_fr^2d\Omega_f$, and the cross-section of the
corresponding reaction by dividing this current by the incident current
density $v_i$:
\begin{equation}
 d\sigma_{fi}=|f_{fi}|^2\frac{p_f}{p_i}d\Omega_f,
 \tag{144.4}
\end{equation}
where $p_i=m_iv_i$, $p_f=m_fv_f$.

In \S~125 the scattering operator $\widehat S$ was introduced as the operator
that transforms an incoming wave into an outgoing wave. With several
channels, this operator has matrix elements for transitions between different
channels. Elements diagonal in the channel indices correspond to elastic
scattering, and off-diagonal elements to the various inelastic processes;
all these elements remain operators in the other variables. They are defined
as follows.

As in \S~125, introduce operators $\widehat f_{ii}$, $\widehat f_{fi}$
associated with the amplitudes $f_{ii}$, $f_{fi}$ by
\begin{equation}
 \widehat S_{fi}=\delta_{fi}+2i\sqrt{k_ik_f}\,\widehat f_{fi}.
 \tag{144.5}
\end{equation}

This definition indeed gives an $S$-matrix that must satisfy unitarity. Write
the entrance-channel wave function as a set of incoming and outgoing waves,
as in \S~125:
\begin{align}
 \psi_i&=F(-\mathbf n')\frac{\exp(-ik_i r)}{r\sqrt{v_i}}
 -(1+2ik_i\widehat f_{ii})F(\mathbf n')\frac{\exp(ik_i r)}{r\sqrt{v_i}}\notag\\
 &=F(-\mathbf n')\frac{\exp(-ik_i r)}{r\sqrt{v_i}}
 -\widehat S_{ii}F(\mathbf n')\frac{\exp(ik_i r)}{r\sqrt{v_i}}.
 \tag{144.6}
\end{align}
For convenience, an extra factor $v_i^{-1/2}$ has been introduced here
relative to (125.3). With the adopted amplitude notation, the wave function
in channel $f$ is
\begin{equation}
 \psi_f=2ik_i\sqrt{m_f/m_i}\,\widehat f_{fi}F(\mathbf n')
 \frac{\exp(ik_f r)}{r\sqrt{v_i}}
 =\widehat S_{fi}F(\mathbf n')\frac{\exp(ik_f r)}{r\sqrt{v_f}}.
 \tag{144.7}
\end{equation}

The incoming-wave current must equal the sum of the outgoing-wave currents in
all channels. This is the evident requirement that the sum of the
probabilities of all

\SourcePage{725}{728}
possible elastic and inelastic processes in a collision be unity. Because of
the factors $\sqrt v$ in the spherical-wave denominators, the velocity drops
out of their current densities. The condition therefore simply requires the
normalization of the incoming wave to equal that of the complete set of
outgoing waves. It remains the unitarity condition for the scattering
operator, now understood also as a matrix in the channel indices. For the
operators $\widehat f_{fi}$ this condition is
\begin{equation}
 \widehat f_{fi}-\widehat f_{if}^{+}
 =2i\sum_n k_n\widehat f_{fn}\widehat f_{in}^{+},
 \tag{144.8}
\end{equation}
analogous to (125.7). The index $+$ denotes complex conjugation and
transposition in every matrix index other than the channel number.

The $S$-matrix is diagonal in states of specified orbital angular momentum
$l$; distinguish the corresponding matrix elements by a superscript $(l)$.
Applying $\widehat f_{ii}$ and $\widehat f_{fi}$ to (125.17) gives the elastic
and inelastic amplitudes
\begin{align}
 f_{ii}&=\frac1{2ik_i}\sum_{l=0}^{\infty}(2l+1)(S_{ii}^{(l)}-1)P_l(\cos\theta),\notag\\
 f_{fi}&=\frac1{2i\sqrt{k_ik_f}}\sum_{l=0}^{\infty}(2l+1)S_{fi}^{(l)}P_l(\cos\theta).
 \tag{144.9}
\end{align}
The corresponding integral cross-sections are
\begin{equation}
 \sigma_{ii}=\frac\pi{k_i^2}\sum_{l=0}^{\infty}(2l+1)|1-S_{ii}^{(l)}|^2,
 \qquad
 \sigma_{fi}=\frac\pi{k_i^2}\sum_{l=0}^{\infty}(2l+1)|S_{fi}^{(l)}|^2.
 \tag{144.10}
\end{equation}
The first agrees with (142.3). The total reaction cross-section $\sigma_r$
from entrance channel $i$ is $\sigma_r=\sum_f'\sigma_{fi}$ over all
$f\ne i$. By unitarity, $\sum_f'|S_{fi}|^2=1-|S_{ii}|^2$, and (142.4) for
$\sigma_r$ is recovered.

Time-reversal symmetry of the scattering process, the reciprocity theorem, is
expressed by
\begin{equation}
 \widehat S_{fi}=\widehat S_{i^*f^*},
 \tag{144.11}
\end{equation}
or equivalently
\begin{equation}
 \widehat f_{fi}=\widehat f_{i^*f^*}.
 \tag{144.12}
\end{equation}

\SourcePage{726}{729}
Here $i^*$ and $f^*$ denote states obtained from $i$ and $f$ by reversing the
particle momenta and spin projections\footnote{For composite particles such
as atoms or atomic nuclei, ``spin'' here means the total intrinsic angular
momentum, formed from both the spins and the orbital angular momenta of the
internal motions of the constituents.}; they are called the
\emph{time-reversed} states of $i$ and $f$. Equations (144.11), (144.12)
generalize (125.11), (125.12) for elastic scattering\footnote{We disregard
here a factor $-1$ that can arise in collisions of particles with spin
(compare (140.11)). This does not, of course, affect the cross-section
relation (144.13).}.

Equation (144.12) gives the following relation between reaction
cross-sections:
\begin{equation}
 \frac{d\sigma_{fi}}{p_f^2d\Omega_f}
 =\frac{d\sigma_{i^*f^*}}{p_i^2d\Omega_{i^*}}.
 \tag{144.13}
\end{equation}
This expresses the \emph{principle of detailed balance}.

As noted in \S~126, when perturbation theory applies, its first approximation
provides, besides the reciprocity theorem, an additional relation between the
amplitudes of the direct and reverse processes in the literal sense,
$i\to f$ and $f\to i$. This property, $f_{fi}=f_{if}^*$, holds in the same
approximation for inelastic processes as well. The cross-sections then obey
\begin{equation}
 \frac{d\sigma_{fi}}{p_f^2d\Omega_f}
 =\frac{d\sigma_{if}}{p_i^2d\Omega_i}.
 \tag{144.14}
\end{equation}

The distinction between $i\to f$ and $i^*\to f^*$ disappears for integral
cross-sections integrated over all directions of $\mathbf p_f$, summed over
the spin directions of the final particles $s_{1f}$, $s_{2f}$, and averaged
over the direction of $\mathbf p_i$ and the spins $s_{1i}$, $s_{2i}$ of the
initial particles. Denote this cross-section by $\overline{\sigma}_{fi}$:
\[
 \overline{\sigma}_{fi}=\frac1{4\pi(2s_{1i}+1)(2s_{2i}+1)}
 \sum_{(m_s)}\int d\sigma_{fi}\,d\Omega_i.
\]
The sum is over the spin projections of all particles. The factor before the
sum and integral occurs because quantities belonging to the initial particles
are averaged rather than summed. Write (144.13) as
\[
 p_i^2d\sigma_{fi}\,d\Omega_{i^*}
 =p_f^2d\sigma_{i^*f^*}\,d\Omega_f
\]
and perform the stated operations. The required relation is
\begin{equation}
 g_ip_i^2\overline{\sigma}_{fi}=g_fp_f^2\overline{\sigma}_{if};
 \tag{144.15}
\end{equation}

\SourcePage{727}{730}
where
\begin{equation}
 g_i=(2s_{1i}+1)(2s_{2i}+1),\qquad
 g_f=(2s_{1f}+1)(2s_{2f}+1),
 \tag{144.16}
\end{equation}
are the numbers of possible spin orientations of the initial and final
particle pairs. These numbers are called the \emph{statistical weights} of
states $i$ and $f$.

Finally, note the following property of $f_{fi}$. We saw in the preceding
section that, as $p_i\to0$, the reaction cross-section varies as
$\sigma_{fi}\sim1/p_i$ when the interaction decreases sufficiently rapidly
at large distances. According to (144.4), this means that
$f_{fi}\to\mathrm{const}$ as $p_i\to0$. By the symmetry (144.12), $f_{fi}$
also tends to a constant as $p_f\to0$. We shall return to this property in
\S~147.
